\section{Cryptographic constructions}

The one-way functions and pseudorandom objects below use the library's bit-string
security experiments and machine-certified efficiency predicates. Security is
asymptotic against uniform strict polynomial-time probabilistic adversaries.
In the one-wayness game, $x$ is uniform among $n$-bit strings; the inverter receives
unary $n$ and $f(x)$ and succeeds with any preimage, even one of a different length.
The function need not preserve length.

\begin{result}
\subsection{Goldreich--Levin hard-core predicate}
For every one-way function $f$, the inner product $\langle x,r\rangle$ over
$\mathbb F_2$ is a hard-core predicate for the padded map $(x,r)\mapsto(f(x),r)$.
The formal construction splits a seed into equal halves and retains an optional
spare bit, covering every seed length. The predicate is polynomial-time computable;
every uniform PPT predictor has negligible advantage over a fair guess.
\citep{GoldreichLevin1989}
\decl{Complexity/Crypto/Computational/GoldreichLevin/HardCore.lean}{159}{Cslib.Crypto.GoldreichLevin.padded_hardCore}
\end{result}

\begin{result}
\subsection{Pseudorandom generators from arbitrary one-way functions}
Every uniform one-way function implies a deterministic polynomial-time
pseudorandom generator expanding each $n$-bit seed to $n+1$ bits, secure against
uniform strict PPT distinguishers. The proof follows Holenstein's
pseudoentropy route to the H\aa stad--Impagliazzo--Levin--Luby theorem.
OWF existence is the sole cryptographic assumption: no entropy estimate or
nonuniform advice is supplied as an extra hypothesis. The result does not claim
optimized construction exponents.\citep{HILL1999,Holenstein2006}
\decl{Complexity/Crypto/Computational/OneWayToPRG.lean}{54}{Cslib.Crypto.OneWay.exists_pseudorandomGenerator}
\end{result}

\begin{result}
\subsection{Polynomial stretch with quantitative resource bounds}
Let $G$ be a one-bit-stretch PRG and let $k:\mathbb N\to\mathbb N$ be positive
and polynomial-time computable with unary input and output. Repeatedly expanding
the retained $n$-bit seed yields a PRG with output length $n+k(n)$, secure against
the same uniform strict PPT distinguishers. In particular, one-bit stretch
suffices for any strictly larger, unary-polynomial-time computable target length.
\citep{Barak2016PRG}

The construction also has a separate machine-level resource certificate.
Suppose $G$ has time/workspace bounds $t_G,s_G$, the count computation has bounds
$t_k,s_k$, and every round's seed $u_i$ satisfies $s_G(u_i)\le Q(u)$ for initial
seed $u$. Writing $n=|u|$ and $k=k(n)$, the certified bounds are
\[
\begin{aligned}
\text{time}&=O\!\left(t_k(n)+n+k+\sum_{i<k}(t_G(u_i)+n+1)+1\right),\\
\text{workspace}&=O\!\left(s_k(n)+Q(u)+n+k+1\right).
\end{aligned}
\]
The exact one-bit output-length invariant is required. The workspace bound
includes the final reversal buffer; a small streaming state alone does not give
an output-length-independent bound for the full construction.
\decl{Complexity/Crypto/Computational/Stretch.lean}{974}{Cslib.Crypto.PRGStretch.pseudorandomGenerator}
\decl{Complexity/Crypto/Computational/Stretch.lean}{429}{Cslib.Crypto.PRGStretch.computableTidilyInTimeAndSpace_stretch}
\end{result}

\begin{result}
\subsection{GGM and adaptive pseudorandom functions}
Every uniform pseudorandom generator gives an efficiently evaluable
pseudorandom function family by polynomial stretching and the GGM tree
construction. Security holds against uniform adaptive strict polynomial-time
oracle adversaries, not just fixed query lists. At parameter $n$, keys, valid
queries and answers have $n$ bits; both games answer malformed queries with the
empty word. Adversaries receive the unary parameter and an empty auxiliary word.
Consequently every uniform
one-way function implies such a PRF. These are conditional security implications,
not assertions that one-way functions exist.\citep{GGM1986,Holenstein2006}
\decl{Complexity/Crypto/Computational/GGM/Security.lean}{97}{Cslib.Crypto.PseudorandomGenerator.exists_pseudorandomFunction}
\decl{Complexity/Crypto/Computational/OneWayToPRF.lean}{46}{Cslib.Crypto.OneWay.exists_pseudorandomFunction}
\end{result}

\begin{result}
\subsection{Generators imply one-way functions}
A secure PRG $G$ whose output length $\ell(n)$ satisfies $\ell(n)\ge2n$
for every seed length $n$ is itself one-way, including against inverters that
return preimages of different lengths. More generally, $G$ is one-way if
$\ell$ is injective and $n\mapsto2^{-(\ell(n)-n)}$ is negligible.
An arbitrary PRG with positive stretch yields a one-way function after
truncation and stretch amplification; the claim in this case is existence of
a possibly different function. Together with HILL, this gives
\[
\text{a one-way function exists}\quad\Longleftrightarrow\quad
\text{a pseudorandom generator exists}.
\]
All security notions are the uniform strict-PPT notions specified above.
Neither side's unconditional existence is asserted.\citep{KatzLindell2020,HILL1999}
\decl{Complexity/Crypto/Computational/PRGToOneWay.lean}{164}{Cslib.Crypto.PseudorandomGenerator.oneWay}
\decl{Complexity/Crypto/Computational/PRGToOneWay.lean}{187}{Cslib.Crypto.PseudorandomGenerator.oneWay_of_injective}
\decl{Complexity/Crypto/Computational/PRGToOneWay.lean}{214}{Cslib.Crypto.exists_oneWay_iff_exists_pseudorandomGenerator}
\end{result}

\begin{result}
\subsection{One-way functions separate P and BPP from NP}
For $f$ computable in polynomial time, the preimage-prefix problem
$\{(1^n,y,u)\mid\exists x\sqsupseteq u,\ |x|=n,\ f(x)=y\}$ is in NP. If it is in P,
bit-by-bit search gives a deterministic polynomial-time inverter; if it is in BPP,
amplification to error $2^{-(L+1)}$ on instances of encoded length $L$ and a union bound over
the $n$ queries of the search give a uniform strict PPT inverter succeeding with probability at
least $1/2$. Hence a uniform one-way function puts its preimage-prefix problem outside P and
BPP, and, by Cook--Levin and closure of P and BPP under polynomial-time reductions,
$\mathrm{SAT}\notin\mathrm{P}$ and $\mathrm{SAT}\notin\mathrm{BPP}$. The existence of
one-way functions is not asserted.\citep{Goldreich2001,KatzLindell2020}
\declat{ad94a1a8bcb2196673cad3ba1fbc37a052ed8532}{Complexity/Crypto/Computational/OneWay/NP.lean}{135}{Cslib.Crypto.preimagePrefix_mem_NP}
\declat{ad94a1a8bcb2196673cad3ba1fbc37a052ed8532}{Complexity/Crypto/Computational/OneWay/NP.lean}{287}{Cslib.Crypto.OneWay.P_ne_NP}
\declat{ad94a1a8bcb2196673cad3ba1fbc37a052ed8532}{Complexity/Crypto/Computational/OneWay/NP.lean}{292}{Cslib.Crypto.OneWay.SAT_notMem_P}
\declat{73cf7ce29398479e7b7fef5c8921b67ae6b7f269}{Complexity/Crypto/Computational/OneWay/BPP.lean}{236}{Cslib.Crypto.OneWay.not_NP_subset_BPP}
\declat{ad94a1a8bcb2196673cad3ba1fbc37a052ed8532}{Complexity/Crypto/Computational/OneWay/BPP.lean}{250}{Cslib.Crypto.OneWay.SAT_notMem_BPP}
\declat{ad94a1a8bcb2196673cad3ba1fbc37a052ed8532}{Complexity/Computability/Complexity/Classes/BPP/Reduction.lean}{62}{Cslib.Complexity.PolyTimeReducible.mem_BPP}
\end{result}

\begin{result}
\subsection{Kabanets--Cai: MCSP in BPP rules out one-way functions}
If MCSP, with instances encoded by the $2^m$-bit truth table followed by the size
threshold in binary and De Morgan circuit size counting all gates, is in BPP, then
no pseudorandom function family secure against uniform adaptive strict oracle-PPT
adversaries exists, hence no uniform one-way function exists; in particular this
holds if MCSP is in P. The hypothesis is MCSP $\in$ BPP at all thresholds; the
distinguisher reads the truth table of the oracle on the $2^m$ words with $n-m$ leading
zeros, $m=(c+2)\lceil\log_2 n\rceil$, and runs the decider only at threshold
$\lfloor 2^m/(5m)\rfloor$: restrictions of a polynomial-time family have circuits of size
at most $c(3n+2)^c+2$ uniformly in the key, while for $m\ge16$ at most $2^{2^m/4}$ of the
$2^{2^m}$ truth tables have circuits with at most $\lfloor 2^m/(5m)\rfloor$ gates. The
advantage is at least $1/4$ for all large $n$.\citep{KabanetsCai2000,GGM1986}
\declat{ad94a1a8bcb2196673cad3ba1fbc37a052ed8532}{Complexity/Crypto/Computational/OneWay/MCSP.lean}{372}{Cslib.Crypto.not_pseudorandomFunction_of_MCSP_mem_BPP}
\declat{ad94a1a8bcb2196673cad3ba1fbc37a052ed8532}{Complexity/Crypto/Computational/OneWay/MCSP.lean}{442}{Cslib.Crypto.not_oneWay_of_MCSP_mem_BPP}
\declat{ad94a1a8bcb2196673cad3ba1fbc37a052ed8532}{Complexity/Crypto/Computational/OneWay/MCSP.lean}{448}{Cslib.Crypto.not_oneWay_of_MCSP_mem_P}
\declat{73cf7ce29398479e7b7fef5c8921b67ae6b7f269}{Complexity/Crypto/Computational/OneWay/MCSP.lean}{296}{Cslib.Crypto.KabanetsCai.isOraclePPT_mcspDistinguisher}
\declat{73cf7ce29398479e7b7fef5c8921b67ae6b7f269}{Complexity/Crypto/Computational/PseudorandomFunction/TruthTable.lean}{173}{Cslib.Crypto.KabanetsCai.exists_complexity_restriction_le}
\end{result}

\begin{result}
\subsection{Gentry's bootstrapping theorem}
Let $S$ be a public-key bit-encryption scheme (\texttt{EncryptionScheme}) equipped with a
bootstrappable augmented-decryption evaluator (\texttt{Bootstrappable}) that homomorphically
evaluates the De Morgan operations (\texttt{Boolean.Op}: Boolean constants \texttt{.const}\,$b$,
unary negation \texttt{.not}, and binary \texttt{.and} and \texttt{.or}) composed with decryption on
pairs of ciphertexts under $\mathsf{sk}_1$ using an encrypted evaluation key
$\mathsf{ek}\leftarrow\mathsf{encryptKey}(\mathsf{pk}_2,\mathsf{sk}_1)$ under a target public key
$\mathsf{pk}_2$, producing ciphertexts under $\mathsf{pk}_2$ of length at most a polynomial bound
$\mathsf{cipherBound}(n)$.
\begin{itemize}
\item \textbf{Leveled FHE without circular security:} For every unary polynomial-time depth bound
  $D:\mathbb N\to\mathbb N$, the acyclic key-chain construction $\mathsf{leveledKeyGen}(D(n),n)$ and
  synchronous level-by-level evaluator $\mathsf{evalLeveledOutput}$ form a compact leveled FHE
  scheme that evaluates every De Morgan circuit $c$ (or formula $F$) of depth at most $D(n)$ in
  deterministic polynomial time, with output ciphertext bit-length at most $\mathsf{cipherBound}(n)$
  independent of $|c|$ and $D(n)$. If $S$ is semantically secure against uniform strict PPT
  distinguishers (\texttt{SemanticSecure}), then the leveled scheme is semantically secure
  (\texttt{LeveledSemanticSecure}) given the full evaluation key chain, proved via a two-tier hybrid
  reduction over the $D(n)$ levels and the $\mathsf{skLength}(n)$ secret-key bits at each level.
\item \textbf{Fixed-key FHE under circular security:} Under circular security
  (\texttt{CircularSecure}, semantic security given a self-encrypted evaluation key
  $\mathsf{ek}\leftarrow\mathsf{encryptKey}(\mathsf{pk},\mathsf{sk})$), the fixed-key scheme
  $(\mathsf{fixedKeyGen},\mathsf{encrypt},\mathsf{evalHomOutput},\mathsf{decrypt})$ homomorphically
  evaluates De Morgan circuits and formulas of arbitrary depth and size in deterministic polynomial
  time with output ciphertext bit-length at most $\mathsf{cipherBound}(n)$.
\end{itemize}
Here \texttt{Bootstrappable} axiomatizes polynomial-time augmented-decryption gate evaluation
(\texttt{evalAugDecrypt}) and its compactness bound (\texttt{cipherBound}) directly as a primitive
interface on $S$, rather than compiling \texttt{S.decrypt} to a circuit and evaluating that circuit
inside a bounded-depth somewhat homomorphic scheme; the theorem also does not construct a concrete
ideal-lattice or LWE-based somewhat homomorphic encryption scheme.\citep{Gentry2009}
\declat{9003d00af791faf5861a7e4cdc233cecc46ab39e}{Complexity/Crypto/Computational/Bootstrapping.lean}{1946}{Cslib.Crypto.EncryptionScheme.leveledFHE_of_bootstrappable}
\declat{9003d00af791faf5861a7e4cdc233cecc46ab39e}{Complexity/Crypto/Computational/Bootstrapping.lean}{1038}{Cslib.Crypto.EncryptionScheme.leveledSemanticSecure_of_semanticSecure}
\declat{9003d00af791faf5861a7e4cdc233cecc46ab39e}{Complexity/Crypto/Computational/Bootstrapping.lean}{1503}{Cslib.Crypto.EncryptionScheme.fixedKeyFHE_of_circularSecure}
\declat{9003d00af791faf5861a7e4cdc233cecc46ab39e}{Complexity/Crypto/Computational/Bootstrapping.lean}{1865}{Cslib.Crypto.EncryptionScheme.decrypt_evalLeveledOutput_circuit}
\declat{9003d00af791faf5861a7e4cdc233cecc46ab39e}{Complexity/Crypto/Computational/Bootstrapping.lean}{1384}{Cslib.Crypto.EncryptionScheme.decrypt_evalHomOutput_circuit}
\end{result}

\begin{result}
\subsection{Digital signatures, Lamport's one-time scheme, and one-way functions}
A uniform digital signature scheme (\texttt{SignatureScheme}) consists of strict PPT algorithms
$\mathsf{keyGen}(1^n)\to(\mathsf{pk},\mathsf{sk})$ and $\mathsf{sign}(\mathsf{sk},m)\to\sigma$ and a
deterministic polynomial-time verifier $\mathsf{verify}(\mathsf{pk},m,\sigma)\in\{\mathsf{true},\mathsf{false}\}$
that accepts every signature in the support of $\mathsf{sign}(\mathsf{sk},m)$ for every $n$-bit
message $m$ and $(\mathsf{pk},\mathsf{sk})$ in the support of $\mathsf{keyGen}(1^n)$.
In the adaptive chosen-message experiment (\texttt{eufCmaGame}, \texttt{EUFCMA}), a uniform strict
oracle-PPT forger (\texttt{IsOraclePPT}) receives $(1^n,\mathsf{pk})$ and adaptive oracle access to
$\mathsf{sign}(\mathsf{sk},\cdot)$ (which logs each queried message) and wins by outputting a pair
$(m^\star,\sigma^\star)$ with $|m^\star|=n$, $m^\star$ not in the query log, and
$\mathsf{verify}(\mathsf{pk},m^\star,\sigma^\star)=\mathsf{true}$. In the one-time experiment
(\texttt{eufOneTimeGame}, \texttt{OneTimeEUFCMA}), the signing oracle answers at most the first
query and returns the empty word without logging on any subsequent query; full \texttt{EUFCMA}
implies \texttt{OneTimeEUFCMA} via a stateful query-limiting wrapper (\texttt{limitToOneQuery}).
\begin{itemize}
\item \textbf{Secure signatures imply one-way functions:} If $S$ is \texttt{OneTimeEUFCMA} (or
  \texttt{EUFCMA}), reindexing a total polynomial-time seeded evaluator
  $\mathsf{seed}\mapsto(n,\mathsf{pk}_n(\mathsf{seed}))$ for $S.\mathsf{keyGen}$ across all input
  lengths via \texttt{GeneratorReindex.parameter} yields a uniform one-way function: any PPT
  inverter induces a zero-query PPT forger that samples a dyadic seed length, recovers a candidate
  seed with the same public key $\mathsf{pk}$, and signs $0^n$ under the corresponding secret key.
\item \textbf{Lamport's one-time signature scheme from any one-way function:} Given any
  polynomial-time function $f:\{0,1\}^*\to\{0,1\}^*$, $\mathsf{lamportScheme}(f)$ samples $2n$
  independent $n$-bit words $(w_0,\dots,w_{2n-1})$, publishes the framed images
  $(f(w_0),\dots,f(w_{2n-1}))$ as $\mathsf{pk}$, and signs $m\in\{0,1\}^n$ by revealing the $n$ words
  $w_{2i+m_i}$ ($0\le i<n$). If $f$ is one-way (\texttt{OneWay}\,$f$), then
  $\mathsf{lamportScheme}(f)$ is \texttt{OneTimeEUFCMA} secure: any successful one-time forgery
  $m^\star\in\{0,1\}^n$ differs from the single queried message (if any) at some bit position
  $i<n$, so the reduction that embeds its challenge $y=f(x^\star)$ at a uniformly chosen dyadic slot
  $j<2n$ and simulates the one-time signing oracle using a dummy word at slot $j$ inverts $f$ with
  probability at least $\mathsf{Adv}(n)/\mathsf{dyadicSize}(2n)$.
\end{itemize}
Consequently, a uniform one-way function exists if and only if a uniform \texttt{OneTimeEUFCMA}
signature scheme exists, and a uniform \texttt{EUFCMA} signature scheme implies a uniform one-way
function. Upgrading one-time signatures to stateless many-time \texttt{EUFCMA} signatures from
arbitrary one-way functions (Rompel's construction of universal one-way hash functions and
GMR/Merkle authentication trees) is not formalized here.\citep{Lamport1979,GoldwasserMicaliRivest1988,Rompel1990}
\decl{Complexity/Crypto/Computational/Signature.lean}{677}{Cslib.Crypto.SignatureScheme.EUFCMA.oneTimeEUFCMA}
\decl{Complexity/Crypto/Computational/OneWayToSignature.lean}{128}{Cslib.Crypto.SignatureScheme.OneTimeEUFCMA.exists_oneWay}
\decl{Complexity/Crypto/Computational/OneWayToSignature.lean}{289}{Cslib.Crypto.SignatureScheme.EUFCMA.exists_oneWay}
\decl{Complexity/Crypto/Computational/OneWayToSignature.lean}{716}{Cslib.Crypto.OneWay.oneTimeEUFCMA_lamport}
\decl{Complexity/Crypto/Computational/OneWayToSignature.lean}{921}{Cslib.Crypto.OneWay.exists_oneTimeSignature}
\decl{Complexity/Crypto/Computational/OneWayToSignature.lean}{927}{Cslib.Crypto.exists_oneWay_iff_exists_oneTimeSignature}
\end{result}

